\section{Introduction}
\label{sec:intro}

Clinical risk scores are used to decide whom to test, treat or monitor, so the
probabilities they output must be meaningful in their own right. A model that ranks
patients well but reports a 40\% risk for a group whose observed event rate is 15\%
can mislead decisions as badly as a model with poor discrimination. Calibration has
therefore been called the ``Achilles heel'' of predictive analytics~\cite{vancalster2019},
and modern reporting and evaluation frameworks ask for discrimination and calibration
together~\cite{steyerberg2010,collins2015,vancalster2016,huang2020}. In high-stakes
settings it is also argued that models should be \emph{inherently} interpretable rather
than explained after the fact~\cite{rudin2019}, which is why point-based scoring
systems remain popular in medicine~\cite{ustun2016}.

Fuzzy rule-based systems~\cite{zadeh1965,mamdani1975,takagi1985} are a natural
candidate for such scores: linguistic rules such as ``IF glucose is High AND BMI is
High'' are readable by clinicians, and genetic algorithms can learn both the rules and
the membership functions from data~\cite{herrera2008,fernandez2019}. Genetic-fuzzy and
neuro-fuzzy classifiers have a long history in medical
applications~\cite{penareyes1999,nauck1999,kahramanli2008,anooj2012,seera2014,pota2017}.
Multi-objective evolutionary fuzzy systems (MOEFSs) made the trade-off between
accuracy and rule-base complexity explicit, returning a Pareto front of models instead
of a single one~\cite{ishibuchi1997,ishibuchi2007,alcala2009,antonelli2011,fazzolari2013,gorzalczany2017}.

Three gaps remain. First, the objectives optimised by genetic-fuzzy medical models are,
almost without exception, accuracy-type criteria (classification rate, error, $F_1$),
sometimes paired with complexity~\cite{fazzolari2013,gacto2011}. None of them is a
proper scoring rule~\cite{gneiting2007}, so nothing in the search rewards
well-calibrated probabilities. Second, the resulting systems are rarely compared with
tuned gradient-boosted decision trees (GBDTs), which remain the strongest general
learners on medium-sized tabular data~\cite{grinsztajn2022,shwartzziv2022}, or with
logistic regression, which in clinical prediction is often as good as more flexible
learners~\cite{christodoulou2019}. Third, comparisons are seldom protected against
model-selection bias and leakage~\cite{cawley2010,kaufman2012}, and design decisions
are often taken on the same data used to report results.

This paper addresses these gaps with a fuzzy risk score that is optimised by NSGA-II~\cite{deb2002}
for three objectives at once: discrimination ($1-$AUROC), calibration (Brier
score~\cite{brier1950}, a strictly proper scoring rule) and interpretability (total
number of rule literals). The contributions are:
\begin{enumerate}
\item A zero-order Takagi--Sugeno \emph{additive fuzzy risk score}: baseline log-odds
plus rule contributions weighted by fuzzy firing strengths. It reads like a point-based
clinical score but uses graded linguistic conditions. Rule weights are estimated by
penalised logistic regression inside the evolutionary loop (hybrid learning), with the
penalty selected per candidate by inner cross-validated log-loss.
\item A three-objective NSGA-II formulation with an external non-dominated archive, an
integer interpretability objective, and a preregistered rule for choosing one
deployable model from the front.
\item A preregistered, leakage-free benchmark on six public clinical datasets. Two
development datasets were used for all design decisions. The protocol, hypotheses and
SHA-256 hashes of the code were locked before four held-out datasets were downloaded.
The benchmark covers 11 models, including tuned XGBoost, LightGBM and CatBoost, post-hoc
Platt and isotonic calibration, and single-objective genetic-fuzzy baselines.
\item An open, tested implementation (\resNFits{} model fits, single-threaded on a
laptop), whose tables and in-text numbers are regenerated automatically from the
stored results.
\end{enumerate}

Section~\ref{sec:related} reviews related work, Section~\ref{sec:method} describes the
model and the optimiser, Section~\ref{sec:protocol} the preregistered protocol,
Section~\ref{sec:results} the results, Section~\ref{sec:rules} the extracted rule bases,
and Section~\ref{sec:discussion} limitations and threats to validity.
